\section{Coherent witnesses from finite approximation}
\label{ana:coherent}

Several packets distinguish an approximate witness at each requested precision
from one coherent name for an exact witness.  That distinction is essential:
an approximate axis of a rotation can change completely as the precision
increases \cite{repoRotation}.  We first isolate a sufficient condition for coherent selection.
It involves only finite search and a supplied quantitative uniqueness bound.

\begin{theorem}[Approximate witnesses with an isolation modulus]
\label{ana:selection}
For each natural number $s$, let $D_s$ be a nonempty, primitively recursively
listed finite set of rational points in a fixed rational box.  Let $E_s(q)$
be a fixed primitive recursive rational residual evaluator, and suppose
its error relative to a
nonnegative represented residual $E(q)$ is at most $2^{-s-2}$.
Suppose that the following two hypotheses have arithmetic proofs:
\begin{enumerate}
\item some member $q$ of $D_s$ has $E(q)\le 2^{-s-1}$;
\item a supplied nondecreasing primitive recursive function $\sigma$ satisfies
\[
 E(x),E(y)\le 2^{-\sigma(k)}
 \quad\Longrightarrow\quad
 \lVert x-y\rVert_\infty\le 2^{-k}
\]
for all rational points under consideration.
\end{enumerate}
There is a fixed primitive recursive sequence $a_s$ with
$E(a_s)\le 2^{-s}$.  Set $\tau(k)=\max(k,\sigma(k))$.
Then $b_k=a_{\tau(k)}$ is a fast Cauchy name:
\[
 k\le l\quad\Longrightarrow\quad
 \lVert b_k-b_l\rVert_\infty\le 2^{-k}.
\]
If the residual has a supplied uniform continuity modulus, the named limit
has residual zero.
\end{theorem}

\begin{proof}
Scan the finite list $D_s$ and retain the first point with the smallest
rational value of $E_s$.  The minimum is found by rational comparison;
neither equality of represented reals nor a choice principle is used.
For a point $q$ supplied by the first hypothesis,
\[
 E(a_s)\le E_s(a_s)+2^{-s-2}
 \le E_s(q)+2^{-s-2}
 \le E(q)+2^{-s-1}\le 2^{-s}.
\]
If $k\le l$, monotonicity of $\tau$ and $\tau(k)\ge\sigma(k)$ make both residuals
$E(b_k),E(b_l)$ at most $2^{-\sigma(k)}$.  The second hypothesis gives the
stated distance estimate.  To verify residual zero at precision $j$, choose
$k$ at least the residual's continuity modulus for precision $j+2$ and
at least $j+2$.  Since $\tau(k)\ge k$, the residual of the limit is bounded
by the residual of $b_k$ plus its continuity error, each at most
$2^{-j-2}$.  This is a finite estimate for each $j$, not an appeal to
abstract completeness.
\end{proof}

The first hypothesis may be supplied by a finite combinatorial existence
proof, or by a net estimate: an approximate solution together with a
continuity modulus and an explicit grid produces such a point.  The second
hypothesis is additional information, and is the part that prevents
incompatible approximations.  Uniqueness without a primitive recursive
isolation modulus is insufficient for this conclusion.  A merely
qualitative compactness argument does not supply either modulus.

For an arithmetic derivation, the loop invariant is that the retained index
is in the scanned prefix and its rational residual estimate is no larger
than every estimate in that prefix.  The conjunction over the prefix is a
bounded Boolean fold.  Fixing $s$ and the list length before induction makes
this a quantifier-free invariant.  The universal real hypotheses are used
only through the finite observation bounds needed for the requested output.

\section{A quantitative implicit branch from a finite certificate}
\label{ana:implicit}

The Puiseux packet proves a simple-root power-series lemma after exact
algebraic preparation \cite{repoPuiseux}.  Its proof uses $H(0,0)=0$ and an exact nonzero
derivative to make low-order coefficients vanish, so that coefficients
stabilize after finitely many substitutions.  The same finite norm method
admits a more general interface: a sufficiently accurate rational centre
and approximate inverse suffice.  The low-order coefficients need not
vanish.  Quantitative contraction replaces exact stabilization.

This extension concerns local branches with an explicit certificate.  It
does not compute the singularity type or decide whether an arbitrary
Cauchy coefficient is zero.  Those tasks remain subject to the packet's
computability obstruction.

\subsection{Finite weighted norms and the certificate}

Use the norm $|u+iv|_1=|u|+|v|$ on Gaussian rationals and their represented
limits.  This norm is submultiplicative, so no square-root operation is
needed.  For a finite polynomial $V(X)=\sum_k v_kX^k$, define
\[
 \lVert V\rVert_R=\sum_k |v_k|_1R^k\qquad (R>0).
\]
The triangle inequality and finite convolution give
\begin{equation}
 \lVert V+W\rVert_R\le\lVert V\rVert_R+\lVert W\rVert_R,
 \qquad
 \lVert VW\rVert_R\le\lVert V\rVert_R\lVert W\rVert_R.
 \label{ana:normrules}
\end{equation}
To prove the second inequality, expand the coefficient of $X^k$ as
$\sum_{a+b=k}v_aw_b$, apply the triangle inequality, and replace the
finite sum over $a+b=k$ by the corresponding part of the product of the
two weighted sums.  Truncation modulo $X^{K+1}$ can only decrease the norm.

Let
\[
 F(X,Z)=\sum_{i,j} f_{ij}X^iZ^j
\]
be a finite polynomial whose coefficients are supplied by fixed primitive
recursive Gaussian-rational Cauchy evaluators.  Choose dyadic
$R=2^{-r}$ and $\rho=2^{-h}$, with $r,h\in\mathbb N$, and rational
majorants $u_{ij}\ge |f_{ij}|_1$.  Require the two rational inequalities
\begin{equation}
 A_0:=\sum_i u_{i0}R^i\le\rho/2,
 \qquad
 Q:=\sum_{i,j\ge1}j u_{ij}R^i\rho^{j-1}\le1/2.
 \label{ana:contract}
\end{equation}
These inequalities are the numerical certificate.  If a coefficient sample
$\widehat f_{ij}$ has certified error at most $\delta$, then
$u_{ij}=|\widehat f_{ij}|_1+\delta$ is an admissible majorant.  Thus the
certificate is checked by finite rational arithmetic.  Name correctness
is a hypothesis of the represented input interface; a finite sample by
itself does not certify an otherwise arbitrary infinite sequence.

\begin{lemma}[Finite polynomial contraction]
\label{ana:contraction}
Under \eqref{ana:contract}, if $\lVert V\rVert_R,\lVert W\rVert_R\le\rho$,
then
\[
 \lVert F(X,V)\rVert_R\le\rho,
 \qquad
 \lVert F(X,V)-F(X,W)\rVert_R\le\tfrac12\lVert V-W\rVert_R.
\]
Both conclusions also hold after truncating every expression modulo
$X^{K+1}$.
\end{lemma}

\begin{proof}
The factorization
$V^j-W^j=(V-W)\sum_{t<j}V^tW^{j-1-t}$ and
\eqref{ana:normrules} give
\[
 \lVert V^j-W^j\rVert_R
 \le j\rho^{j-1}\lVert V-W\rVert_R.
\]
Multiply by $|f_{ij}|_1R^i$ and sum to obtain the contraction bound.
Taking $W=0$ gives
$\lVert F(X,V)\rVert_R\le A_0+Q\lVert V\rVert_R\le\rho$.
For truncated arrays the same finite convolution proof applies: discarding
terms of degree greater than $K$ cannot increase a weighted sum.
\end{proof}

\subsection{The coefficient construction and all error bounds}

Define finite polynomials
\[
 Z_0=0,\qquad Z_{n+1}=F(X,Z_n).
\]
When only coefficients through $K$ are needed, perform the recursion in
the truncated polynomial ring modulo $X^{K+1}$.  Truncation is a ring
homomorphism, so this gives exactly the same first $K+1$ coefficients.

\begin{theorem}[Certified power-series branch]
\label{ana:series}
Under \eqref{ana:contract}, the following statements have finite arithmetic
proofs:
\begin{enumerate}
\item for $m\ge n$,
\begin{equation}
 \lVert Z_n\rVert_R\le\rho,
 \qquad
 \lVert Z_m-Z_n\rVert_R\le\rho\,2^{-n};
 \label{ana:iterate}
\end{equation}
\item each coefficient has a fixed primitive recursive Cauchy name $c_k$;
when $[X^k]Z_{s+rk+3}$ is approximated with error at most $2^{-s-3}$,
the resulting rational value is a valid stage-$s$ value of that name;
\item for every finite $K$,
\begin{equation}
 \sum_{k=0}^K |c_k|_1R^k\le\rho;
 \label{ana:coefficientbound}
\end{equation}
\item the formal series $C(X)=\sum_{k\ge0}c_kX^k$ satisfies
$C=F(X,C)$ coefficient by coefficient;
\item for $|x|_1\le R/2$ and $M>K$,
\begin{equation}
 \left|\sum_{k=K+1}^M c_kx^k\right|_1
 \le\rho\,2^{-K-1}.
 \label{ana:tail}
\end{equation}
In particular, the coefficient names specify a convergent branch with an
explicit primitive recursive tail modulus.  Among series satisfying the
finite weighted bounds \eqref{ana:coefficientbound}, it is the unique
fixed point.  At each $|x|_1\le R/2$, its value is the unique fixed point
in the ball $|z|_1\le\rho$.
\end{enumerate}
\end{theorem}

\begin{proof}
The norm bound is induction using Lemma~\ref{ana:contraction}.
The first difference has norm at most $A_0\le\rho/2$; applying the
contraction bound repeatedly gives
\[
 \lVert Z_{t+1}-Z_t\rVert_R\le (\rho/2)2^{-t}.
\]
Finite telescoping from $t=n$ to $m-1$ and the finite geometric-sum
identity give \eqref{ana:iterate}.  Consequently
\[
 |[X^k](Z_m-Z_n)|_1\le\rho R^{-k}2^{-n}.
\]
Since $R^{-k}=2^{rk}$ and $\rho\le1$, the stage
$n=s+rk+3$ makes the right side at most $2^{-s-3}$.
If $s\le t$, the two chosen polynomial stages are ordered.  Adding their
two rational evaluation errors yields a distance at most
$3\cdot2^{-s-3}<2^{-s}$.  This proves the fast Cauchy condition directly
for the declared name; no prior limit is used to define it.

For \eqref{ana:coefficientbound}, fix $K$ and a precision $s$.
Approximate all $K+1$ coefficient names to precision
$p=s+K+3$ and compare them to one common stage
$N=p+rK+3$.  The coefficient estimate bounds the comparison between a
stage-$p$ name value and the common polynomial stage by $2^{-p-2}$
per coefficient.  Adding the name error $2^{-p}$ gives at most
$(5/4)2^{-p}$ per coefficient.  Because $R\le1$ and
$K+1\le2^{K+1}$, the total weighted error is at most
$(5/16)2^{-s}<2^{-s}$.  The common polynomial stage has norm at most $\rho$.
This proves the rational meaning of
$\sum_{k\le K}|c_k|_1R^k\le\rho$ at every precision.  The argument is a
finite estimate with an arbitrary parameter $s$, rather than a
completeness argument about a space of all series.

Fix an output degree $K$.  Computation of the first $K+1$ coefficients
of $F(X,V)$ uses only the first $K+1$ coefficients of $V$, by finite
convolution.  Compare the truncation $C^{\le K}$ with $Z_n^{\le K}$.
The preceding coefficient estimates give convergence of their finite
weighted difference, with a primitive recursive modulus.  The contraction
lemma bounds the difference of their images.  The identity
$Z_{n+1}=F(X,Z_n)$ then makes every coefficient of $C-F(X,C)$ arbitrarily
small.  This is exactly coefficientwise equality in the name semantics.

The tail estimate is a finite sum:
\[
 \sum_{k=K+1}^M|c_k|_1|x|_1^k
 \le 2^{-K-1}\sum_{k=K+1}^M|c_k|_1R^k
 \le\rho\,2^{-K-1}.
\]
It defines the evaluated branch by its rational partial sums and an
explicit error budget for finitely many coefficient evaluations.
The finite partial-sum bounds also give $|C(x)|_1\le\rho$.
To verify the pointwise fixed-point identity without presupposing that
evaluation commutes with an infinite substitution, put
$G_K=C_K-F(X,C_K)$, where $C_K=\sum_{k\le K}c_kX^k$.
Coefficientwise equality makes its coefficients through $K$ vanish.
Invariance gives $\lVert G_K\rVert_R\le2\rho$, hence
$|G_K(x)|_1\le2\rho\,2^{-K-1}$ when $|x|_1\le R/2$.
The tail estimate and scalar contraction now imply
\[
 |C(x)-F(x,C(x))|_1
 \le\left(\frac32+2\right)\rho\,2^{-K-1}
 =\frac72\rho\,2^{-K-1}.
\]
Taking $K=j+3$ makes this less than $2^{-j}$.
This proves the pointwise identity by explicit finite residual bounds.
For uniqueness of formal solutions, truncate two candidate fixed points
through $K$.  The contraction estimate in the truncated ring gives
$d\le d/2$ for their weighted distance $d$, hence $d=0$.
This holds for every $K$.  The same factorization gives
$|z-w|_1\le|z-w|_1/2$ for pointwise fixed points in the specified ball.
Both are quantitative equalities; neither uses decidable real equality.
\end{proof}

\subsection{Application to arbitrary represented polynomial coefficients}

Let $P(X,Y)$ have arbitrary fixed primitive recursive Cauchy coefficient
names.  Supply rational Gaussian numbers $a,v$, with $v\ne0$, and form
\[
 F(X,Z)=Z-vP(X,a+Z).
\]
If a certificate \eqref{ana:contract} passes, the theorem constructs a
branch
\[
 Y(X)=a+C(X),\qquad P(X,Y(X))=0.
\]
There is no assumed exact root at $X=0$ and no exact test of any input
coefficient.  A convenient approximate inverse $v$ makes the coefficient
of $Z$ small, and a sufficiently accurate centre $a$ makes the constant
defect small.  The two certificate inequalities state exactly how small
they must be.  The reciprocal of $v$ is finite rational data.

If $C_K=\sum_{k\le K}c_kX^k$, the tail bound and contraction estimate give
the useful residual certificate
\begin{equation}
 |P(x,a+C_K(x))|_1
 \le |v^{-1}|_1\,\frac32\rho\,2^{-K-1}
 \qquad(|x|_1\le R/2).
 \label{ana:presidual}
\end{equation}
Indeed, $|C_K(x)|_1\le\rho$ by the finite weighted bound, and
\[
 |C_K-F(x,C_K)|_1
 \le |C_K-C|_1+|F(x,C)-F(x,C_K)|_1
 \le\tfrac32|C_K-C|_1.
\]
The equality $vP=C_K-F(x,C_K)$ gives \eqref{ana:presidual}.

For a concrete example, let $\alpha$ be any represented real with
$|\alpha|\le1/16$, and put
\[
 P(X,Y)=Y^2-(1+\alpha+X),\qquad a=1,\quad v=1/2.
\]
Then
$F(X,Z)=\alpha/2+X/2-Z^2/2$.
With $R=1/8$ and $\rho=1/4$,
\[
 A_0\le1/32+1/16=3/32<\rho/2,
 \qquad Q=\rho=1/4<1/2.
\]
The construction gives the branch near $Y=1$ uniformly in the name of
$\alpha$, including $\alpha=0$; it never decides that equality.
The constant coefficient is computed by contraction and need not have an
exact algebraic presentation.  This example also permits transcendental
$\alpha$.

The Puiseux packet's singular example
$Y^2-X^2-\alpha$ remains outside any uniform regular-chart assertion at
$(X,Y,\alpha)=(0,0,0)$: there the required derivative inverse disappears.
If a particular ramified substitution and finite chart are supplied and
verified, our theorem applies to its regularized polynomial as well.
It does not discover such a chart uniformly from arbitrary coefficient
names.  Thus it extends the usable coefficient representation on certified
regular charts while preserving the exact boundary established by the
existing obstruction \cite{repoObstruction}.

\subsection{Arithmetic audit and formalization obligations}

For rational coefficient arrays every displayed norm, certificate test,
truncated convolution, and iteration is primitive recursive.  The core
induction invariant is the conjunction of the norm bound and geometric
increment bound at the current stage.  Rational denominators are cleared
by their displayed positive products.  Each finite convolution inequality
uses an induction on a bounded array length.  Multiple invariants are
combined into one Boolean equality, so quantifier-free induction suffices.

For represented coefficients, first fix the requested output degree,
iteration horizon, and error tolerance.  The reduction to exact rational
inputs must preserve the certificate, not only approximate its hypotheses.
For a rational sample $q$ and a supplied majorant $u\ge0$, put
\[
 \pi_u(q)=
 \begin{cases}
 q,&|q|_1\le u,\\
 (u/|q|_1)q,&|q|_1>u.
 \end{cases}
\]
The second denominator is positive in its rational branch.  If
$|f|_1\le u$ and $|q-f|_1\le\delta$, then
$|q|_1-u\le\delta$ whenever clipping occurs.  Thus
$|\pi_u(q)-f|_1\le2\delta$, while $|\pi_u(q)|_1\le u$ in every case.
Clip each sample to its majorant.  The resulting rational polynomial
satisfies exactly the original certificate, so the finite rational proof
applies to every such approximation.

When contraction is applied to arbitrary represented finite arrays
$V,W$ with weighted norms at most $\rho$, their rational samples must
also preserve that ball.  Clip each whole sampled array $v$ radially:
retain it if $\lVert v\rVert_R\le\rho$, and otherwise replace it by
$(\rho/\lVert v\rVert_R)v$.  The same calculation in the finite weighted
norm shows that this has norm at most $\rho$ and increases an approximation
error by at most a factor of two.  The estimate
$\bigl||z|_1-|w|_1\bigr|\le|z-w|_1$ and a finite weighted sum transfer
norm comparisons.  Thus all rational hypotheses needed for contraction,
including in the uniqueness proof, hold exactly before transfer.

Now use the arithmetic-circuit error lemma, Lemma~\ref{core:gate}, to
transfer the desired finite output inequalities.  A concrete primitive
recursive choice is obtained by executing the truncated iteration on
formal expression trees: each array addition, product, and coefficient
copy appends its fixed finite arithmetic circuit, so its gate count $N$
is obtained by a bounded count.  Upper bounds for intermediate magnitudes
are computed by the same circuit with addition of nonnegative bounds for
$+$ and $-$ and multiplication for $\cdot$; take an integer $B\ge1$
larger than all these finitely many bounds.  With $L=2B+3$, sample input
coefficients so that, after clipping, their error is at most
$\delta=2^{-s}/L^{N+1}$.  For example an additional factor of four in the
input dyadic precision covers clipping and rounding.  The gate lemma then
bounds the final error by $2^{-s}$.  All $N,B,L$ and precisions are
primitive recursive functions of the fixed finite program, iteration
count, degree cutoff, and input majorants.  This proves each requested
name inequality from rational finite observations.  No induction over
an infinite-series assertion or an unbounded universal real predicate is
performed.

The accompanying \texttt{SeriesTargets.lean} defines truncated rational
arrays, convolution, iteration, weighted norms, and the principal finite
target propositions.  Completion should prove, in order: convolution
submultiplicativity; the finite power-difference identity; contraction and
invariance; finite geometric telescoping; coefficient extraction bounds;
the explicit Cauchy-name construction; finite substitution continuity;
and the output formulas for the encoded arithmetic.
The first four can be proved over rational arrays before any represented
analysis is introduced.  A proof in Lean of those array lemmas is a useful
prototype; constructing the corresponding object-theory derivations
remains a separate required step.  No claim of a completed iPRA Lean
formalization is made by a record of analytic assumptions.
